\documentclass[10pt,conference]{IEEEtran}

\usepackage[utf8]{inputenc}
\usepackage[T1]{fontenc}
\usepackage[spanish,es-noquoting]{babel}
\usepackage{amsmath,amssymb}
\usepackage{graphicx}
\usepackage{booktabs}
\usepackage{multirow}
\usepackage{cite}
\usepackage{url}
\usepackage{hyperref}
\usepackage{tabularx}
\usepackage{siunitx}

\hypersetup{
    colorlinks=true,
    linkcolor=blue,
    citecolor=blue,
    urlcolor=blue,
    pdftitle={Caracterizacion objetiva y comparativa de emulaciones de cinta magnetica},
}

\begin{document}

\title{Caracterización Objetiva y Comparativa de Emulaciones de Cinta Magnética: \\
Protocolos de Calibración, Mantenimiento de Respuesta frente a Ganancia y Compresión No Lineal por Frecuencia}

\author{
\IEEEauthorblockN{Alioth Ponce Aranda}
\IEEEauthorblockA{APA Audio Labs\\
Investigación en Acústica y Procesamiento Digital de Señales (DSP)\\
Correo electrónico: \texttt{aliothponce@gmail.com}}
}

\maketitle

\begin{abstract}
Se establece un protocolo de medición empírico y objetivamente reproducible para evaluar y
comparar emulaciones de cinta magnética de distribución gratuita. Se analiza el plugin
\emph{Iron Deck} (M~Media Audio) frente a un procesador de referencia representativo de
una firma internacionalmente establecida en hardware de grabación analógico.
La investigación examina la divergencia de estándares de ecualización entre ambos dispositivos
—coincidiendo únicamente en NAB, mientras que Iron Deck ofrece IEC/AES y la referencia expone
CCIR— y evalúa los protocolos de calibración físicos (cintas patrón MRL/STL, flujos magnéticos
de $200$, $250$, $355$ y $510\,\text{nWb/m}$) y su correspondencia en el dominio digital.
Mediante scripts en Python se realizan tres análisis exclusivos: (1)~estabilidad de calibración
espectral bajo excursión de ganancia (\emph{drive sweep}); (2)~curva de transferencia estática
Input$\to$Output a múltiples frecuencias; y (3)~compresión no lineal dependiente de la frecuencia.
Los resultados muestran que Iron Deck iguala la integridad digital del dispositivo de referencia
en estabilidad de frecuencia y anti-aliasing, y ofrece ventajas de versatilidad en la selección
de estándar de ecualización y velocidad de transporte.
Este estudio sienta un precedente metodológico para la comunidad de ingeniería de audio.
\end{abstract}

\begin{IEEEkeywords}
emulación de cinta magnética, cinta patrón MRL, ecualización NAB/IEC/CCIR, respuesta en
frecuencia, distorsión armónica total (THD), compresión no lineal por frecuencia, aliasing,
audio DSP, análisis objetivo.
\end{IEEEkeywords}

% ===========================================================================
\section{Introducción}
% ===========================================================================

En el ámbito del procesamiento de señales de audio, las emulaciones analógicas desempeñan un
rol crítico al impartir cohesión dinámica y coloración armónica características del medio
magnético. Históricamente, la apreciación de estos procesadores ha estado supeditada al
prestigio comercial de la marca desarrolladora o al nombre de la máquina física modelada,
relegando el análisis objetivo a un segundo plano.

La presente investigación aborda una comparativa rigurosa entre dos emulaciones de cinta
magnética de distribución gratuita pero de orígenes opuestos: \emph{Iron Deck} (M~Media
Audio), un plugin de un desarrollador independiente con control paramétrico de velocidad y
curva de ecualización; y un plugin de referencia distribuido por una firma líder en hardware
de grabación analógico, aquí tratado de forma anónima por razones de independencia técnica.

El objetivo central no es determinar subjetivamente cuál dispositivo produce una coloración
sonora más apreciada, sino construir un precedente empírico basado en métricas DSP
reproducibles: respuesta en frecuencia y fase, distorsión armónica total (THD), compresión
no lineal dependiente de la frecuencia, deriva de calibración espectral e inmunidad al
aliasing. Todas las mediciones se realizaron programáticamente para garantizar su
reproducibilidad independiente de criterio perceptual.

% ===========================================================================
\section{Fundamentos Físicos y Estándares de Calibración}
% ===========================================================================

\subsection{Pérdidas Magnéticas y Geometría de Cabezas}

La física del transporte de cinta impone restricciones dependientes de la velocidad $v$ y
la longitud de onda magnética $\lambda = v/f$. A bajas velocidades ($7.5\,\text{ips}$),
la longitud de onda de los agudos ($10\,\text{kHz}$) se reduce drásticamente, incrementando
las pérdidas por entredicho magnético (\emph{gap loss}):
\begin{equation}
L_g(f) = 20\log_{10}\!\left|\frac{\sin(\pi g/\lambda)}{\pi g/\lambda}\right| \quad (\text{dB})
\end{equation}
donde $g$ es el ancho físico del entredicho de la cabeza reproductora. La pérdida de
separación de Wallace \cite{wallace}, $L_d(f)=54.6\,(d/\lambda)\,\text{dB}$, atenúa los
agudos cuando la cinta se despega de la cabeza. En el rango grave, la geometría de los polos
magnéticos introduce interferencia constructiva/destructiva conocida como \emph{head bump},
característica en el rango $50$–$100\,\text{Hz}$.

\subsection{Estándares de Ecualización: Coincidencia en NAB y Divergencia IEC/CCIR}

Un hallazgo fundamental de este estudio es la divergencia de estándares de ecualización entre
los dos dispositivos evaluados:

\begin{itemize}
    \item \textbf{NAB} (National Association of Broadcasters): única norma presente en
    \emph{ambos} plugins. Utiliza constantes de tiempo $\tau_1=50\,\mu\text{s}$
    ($f_1\approx 3183\,\text{Hz}$) y $\tau_2=3180\,\mu\text{s}$ ($f_2\approx 50\,\text{Hz}$),
    con realce en agudos y énfasis en graves (norma americana).

    \item \textbf{IEC~1\,/\,CCIR} (norma europea, $15\,\text{ips}$): disponible únicamente
    en el dispositivo de referencia. Constante $\tau_1=35\,\mu\text{s}$
    ($f_1\approx 4547\,\text{Hz}$), sin énfasis de graves, produciendo un perfil más plano
    y transparente en la banda media-alta.

    \item \textbf{IEC~2\,/\,AES} ($30\,\text{ips}$): disponible únicamente en Iron Deck.
    Constante $\tau_1=17.5\,\mu\text{s}$ ($f_1\approx 9095\,\text{Hz}$), diseñada para
    máxima respuesta extendida en alta frecuencia y menor piso de ruido a alta velocidad.
\end{itemize}

La inclusión de CCIR en la referencia refleja el modelado de equipos de estudio europeos
(p.\,ej.\ Studer, Telefunken), mientras que la oferta de IEC/AES en Iron Deck apunta a
flujos de trabajo de alta resolución a $30\,\text{ips}$. Esta divergencia de diseño es
en sí misma un objeto de estudio técnico.

\subsection{Protocolos Físicos de Calibración y su Mapeo Digital}

En máquinas de cinta físicas, la alineación de estudio requiere cintas patrón certificadas
(MRL/STL) \cite{mrl} para ajustar:
\begin{enumerate}
    \item \textbf{Azimut}: perpendicularidad de la cabeza reproductora mediante tono de
    $10$–$15\,\text{kHz}$.

    \item \textbf{Flujo operativo}: calibración de $0\,\text{VU}\equiv+4\,\text{dBu}
    \equiv -18\,\text{dBFS}$ con tono de $1\,\text{kHz}$ a flujos de $200$, $250$, $355$
    o $510\,\text{nWb/m}$.

    \item \textbf{Ajuste de bias}: corriente de polarización de $100$–$200\,\text{kHz}$
    ajustada por \emph{overbias} ($-1.5\,\text{dB}$ a $10\,\text{kHz}$) para optimizar el
    balance entre distorsión y ruido.
\end{enumerate}

En el dominio digital, este protocolo se evalúa verificando que el procesado no introduzca
derivas de frecuencia ni desfases erráticos ante excursiones de ganancia, lo que constituye
el objeto de la Sección~\ref{sec:calibracion} de este trabajo.

% ===========================================================================
\section{Metodología de Análisis DSP Programático}
% ===========================================================================

Las mediciones se realizaron programáticamente en Python~3 (\texttt{scipy.signal},
\texttt{numpy}) a $48\,\text{kHz}$\,/\,32-bit float, sobre señales generadas y procesadas
por cada plugin en tiempo real bajo las condiciones exactas de cada experimento:

\begin{itemize}
    \item \textbf{Respuesta en frecuencia}: barrido senoidal logarítmico
    ($1\,\text{Hz}$–$24\,\text{kHz}$). La función de transferencia $H(f)$ se estimó mediante
    la relación de densidades espectrales cruzadas (CSD):
    \begin{equation}
        H(f) = \frac{S_{xy}(f)}{S_{xx}(f)}
    \end{equation}
    con suavizado Savitzky-Golay ($w=151$, $p=2$) para eliminar ruido espectral
    sin distorsionar las pendientes de EQ.

    \item \textbf{Transferencia estática y THD}: tonos puros a $100\,\text{Hz}$,
    $1\,\text{kHz}$ y $5\,\text{kHz}$, evaluados a niveles de entrada de $-30$, $-18$,
    $-12$, $-6$ y $0\,\text{dBFS}$ para caracterizar compresión y distorsión armónica
    ($H_2$–$H_{19}$).

    \item \textbf{Estabilidad de calibración}: tonos de referencia a $3000\,\text{Hz}$
    (tono patrón NAB) y $3150\,\text{Hz}$ (DIN/IEC), medidos con Welch ($N_{\text{FFT}}
    = 4f_s$) para resolución sub-Hz.

    \item \textbf{Distribución espectral}: ruido rosa procesado y segmentado en
    10 bandas de octava centradas en $31.5$–$16000\,\text{Hz}$.
\end{itemize}

Los archivos de entrada siguen la convención de nombres:
\begin{center}
\texttt{\{Dispositivo\}-\{IPS\}\_IPS-EQ\_\{curva\}-\{P1\}\_\{V1\}-\{P2\}\_\{V2\}.wav}
\end{center}
que permite extracción automática de metadata (velocidad, estándar EQ, nivel de drive)
mediante expresión regular, sin asumir una rejilla fija de parámetros.

% ===========================================================================
\section{Resultados y Análisis Experimental}
% ===========================================================================

\subsection{Retrato Tonal Nominal: Iron Deck}
\label{sec:retrato_iron}

La Fig.\,\ref{fig:01a_retrato_iron} presenta la comparación directa de la respuesta espectral
$H(f)$ a nivel nominal ($-12\,\text{dBFS}$) de Iron Deck entre los modos NAB ($15\,\text{ips}$)
e IEC ($30\,\text{ips}$). La curva NAB exhibe un \emph{head bump} de aproximadamente
$+0.8\,\text{dB}$ en $75\,\text{Hz}$, característico del realce de graves NAB, con caída
progresiva por encima de $10\,\text{kHz}$. La curva IEC a $30\,\text{ips}$ extiende la
respuesta lineal hasta $22\,\text{kHz}$, con un comportamiento considerablemente más plano
en la banda media-alta. El panel diferencial $\Delta H(f)$ aísla cuantitativamente la
ganancia o pérdida por banda al conmutar entre estándares.

\begin{figure}[htbp]
\centering
\includegraphics[width=0.48\textwidth]{resultados_analisis/01a_Retrato_Tonal_IronDeck.png}
\caption{Retrato tonal nominal de Iron Deck. Comparación entre curva NAB ($15\,\text{ips}$)
e IEC ($30\,\text{ips}$) a nivel de entrada $-12\,\text{dBFS}$. El panel inferior muestra
la diferencia espectral $\Delta H(f)$ (IEC$-$NAB) en dB.}
\label{fig:01a_retrato_iron}
\end{figure}

\subsection{Retrato Tonal Nominal: Dispositivo de Referencia}
\label{sec:retrato_ref}

La Fig.\,\ref{fig:01b_retrato_ref} presenta el retrato tonal del dispositivo de referencia
comparando NAB ($15\,\text{ips}$) y CCIR ($15\,\text{ips}$). La curva CCIR muestra atenuación
progresiva por encima de $12\,\text{kHz}$ y ausencia de realce de graves, rasgos
característicos de la norma europea. La diferencia $\Delta H(f)$ respecto a NAB supera
$-3\,\text{dB}$ en las bandas de alta frecuencia, lo que cuantifica el impacto tonal de la
elección de estándar de reproducción.

\begin{figure}[htbp]
\centering
\includegraphics[width=0.48\textwidth]{resultados_analisis/01b_Retrato_Tonal_Referencia.png}
\caption{Retrato tonal nominal del dispositivo de referencia. Comparación entre NAB y
CCIR (ambas a $15\,\text{ips}$, $-12\,\text{dBFS}$). El panel inferior muestra
$\Delta H(f)$ (CCIR$-$NAB) en dB.}
\label{fig:01b_retrato_ref}
\end{figure}

Las Figs.\,\ref{fig:02_velocidad} y \ref{fig:03_eq} detallan el efecto de la velocidad de
transporte ($7.5$, $15$, $30\,\text{ips}$) y el de la curva de ecualización (NAB vs.\ IEC)
en Iron Deck, respectivamente.

\begin{figure}[htbp]
\centering
\includegraphics[width=0.48\textwidth]{resultados_analisis/02_Efecto_Velocidad_Cinta_IronDeck.png}
\caption{Efecto de la velocidad de transporte en Iron Deck ($7.5$, $15$ y $30\,\text{ips}$)
bajo ecualización NAB, a $-12\,\text{dBFS}$ de entrada. A menor velocidad, se acentúan
las pérdidas por \emph{gap loss} en alta frecuencia y el \emph{head bump} en graves.}
\label{fig:02_velocidad}
\end{figure}

\begin{figure}[htbp]
\centering
\includegraphics[width=0.48\textwidth]{resultados_analisis/03_Efecto_Curva_EQ_IronDeck.png}
\caption{Comparación entre ecualizaciones NAB e IEC en Iron Deck a $30\,\text{ips}$.
La norma IEC/AES produce una respuesta más extendida en alta frecuencia y sin el énfasis
de graves propio de NAB.}
\label{fig:03_eq}
\end{figure}

\subsection{Análisis Exclusivo: Mantenimiento de la Calibración frente a la Ganancia}
\label{sec:calibracion}

Para determinar si el control de saturación (\emph{drive}) altera la calibración frecuencial
del plugin, se define la deriva tonal relativa:
\begin{equation}
\Delta H(f,\,D) = |H(f,D)|_{\text{dB}} - |H(f,D_{\text{nom}})|_{\text{dB}}
\end{equation}
donde $D_{\text{nom}}$ es el punto de operación nominal (control REC~=~5 para Iron Deck,
Input~=~0~dB para la referencia). La Fig.\,\ref{fig:09_calibracion} presenta esta deriva
para ambos dispositivos a $15\,\text{ips}$, NAB.

\begin{figure}[htbp]
\centering
\includegraphics[width=0.48\textwidth]{resultados_analisis/09_Mantenimiento_Calibracion_vs_Ganancia.png}
\caption{Mantenimiento de la calibración espectral en función del nivel de \emph{drive}
(panel superior: $H(f)$ por punto de drive; panel inferior: deriva $\Delta H(f)$ respecto
al punto nominal). Izquierda: Iron Deck (REC 0\,/\,5\,/\,10). Derecha: Referencia
(Input $-12$\,/\,$0$\,/\,$+24$~dB).}
\label{fig:09_calibracion}
\end{figure}

Los resultados confirman que Iron Deck mantiene la calibración de respuesta espectral dentro
de un margen $|\Delta H(f)| < 0.15\,\text{dB}$ para cualquier punto de drive en la banda
$20\,\text{Hz}$–$20\,\text{kHz}$, añadiendo saturación armónica sin alterar el balance
espectral. El dispositivo de referencia exhibe una deriva mayor a $+24\,\text{dB}$ de input,
con comprensión en alta frecuencia de hasta $\sim$$0.5\,\text{dB}$, resultado del límite
de la curva de magnetización a excitación extrema.

\subsection{Análisis Exclusivo: Curva de Transferencia Estática}
\label{sec:transferencia}

La Tabla\,\ref{tab:transferencia} sintetiza la \emph{compresión relativa de nivel}
($\Delta_{\text{REC}} = \overline{\Delta_C}(\text{REC}_{\text{max}}) -
\overline{\Delta_C}(\text{REC}_{\text{min}})$) medida a $1\,\text{kHz}$ y $100\,\text{Hz}$
para cada punto de drive evaluado. $\Delta_C = \text{Out}_{\text{RMS}} -
\text{In}_{\text{RMS}}$ (en dB) cuantifica la desviación de la ganancia de inserción
respecto a la relación lineal 1:1; $\Delta_{\text{REC}}$ es la excursión total de compresión
al recorrer el rango completo del control de drive.

\begin{table}[htbp]
\caption{Excursión de compresión al recorrer el rango completo de drive.
$\Delta_C = \text{Out}_\text{RMS} - \text{In}_\text{RMS}$ (dB) medida a entrada
$-3\,\text{dBFS}$. $\Delta_\text{REC}$ es la diferencia de $\Delta_C$ entre el
drive máximo y mínimo evaluados ($15\,\text{ips}$, NAB).}
\label{tab:transferencia}
\centering
\small
\begin{tabular}{llcrr}
\toprule
Dispositivo & Drive & $f$ & $\Delta_C$ (dB) & $\Delta_\text{REC}$ (dB) \\
\midrule
Iron Deck & REC 0  & 1\,kHz & ref.\,$\Delta_C^0$      & \multirow{3}{*}{$-19.0$} \\
Iron Deck & REC 5  & 1\,kHz & $\Delta_C^0 - 3.67$  & \\
Iron Deck & REC 10 & 1\,kHz & $\Delta_C^0 - 18.99$ & \\
\midrule
Iron Deck & REC 0  & 100\,Hz & ref.\,$\Delta_C^0$      & \multirow{3}{*}{$-19.1$} \\
Iron Deck & REC 5  & 100\,Hz & $\Delta_C^0 - 3.74$  & \\
Iron Deck & REC 10 & 100\,Hz & $\Delta_C^0 - 19.08$ & \\
\midrule
Referencia & Input $-12$\,dB & 1\,kHz  & ref.\,$\Delta_C^0$      & \multirow{3}{*}{$-21.8$} \\
Referencia & Input $+0$\,dB  & 1\,kHz  & $\Delta_C^0 - 0.03$  & \\
Referencia & Input $+24$\,dB & 1\,kHz  & $\Delta_C^0 - 21.85$ & \\
\midrule
Referencia & Input $-12$\,dB & 100\,Hz & ref.\,$\Delta_C^0$      & \multirow{3}{*}{$-22.7$} \\
Referencia & Input $+0$\,dB  & 100\,Hz & $\Delta_C^0 - 0.47$  & \\
Referencia & Input $+24$\,dB & 100\,Hz & $\Delta_C^0 - 22.66$ & \\
\bottomrule
\end{tabular}
\end{table}

La Tabla\,\ref{tab:transferencia} muestra que la excursión de compresión al recorrer el
rango completo de drive es de $\approx 19\,\text{dB}$ en Iron Deck y de
$\approx 21$--$23\,\text{dB}$ en la referencia, con comportamiento similar a $1\,\text{kHz}$
y $100\,\text{Hz}$.

\subsection{Análisis Exclusivo: Compresión No Lineal por Frecuencia}
\label{sec:compresion_frec}

La Tabla\,\ref{tab:compresion} presenta el análisis de la variación de $\Delta_C$ al
recorrer el rango de nivel de entrada ($-33$ a $-3\,\text{dBFS}$) a ajuste nominal de
drive, para tres bandas de frecuencia. La columna ``Excursión'' reporta
$\delta = \Delta_C^{\text{máx}} - \Delta_C^{\text{mín}}$, indicando en cuántos dB
varía la compresión efectiva al barrer todo el rango dinámico: una excursión mayor implica
una curva de transferencia más no lineal.

\begin{table}[htbp]
\caption{Variación de la compresión neta $\Delta_C = \text{Out}_\text{RMS} -
\text{In}_\text{RMS}$ al recorrer el rango de entrada $-33$ a $-3\,\text{dBFS}$,
a ajuste nominal de drive ($15\,\text{ips}$, NAB). La excursión
$\delta = \Delta_C^{\text{máx}} - \Delta_C^{\text{mín}}$ indica el grado
de no linealidad de la curva de transferencia en cada banda.}
\label{tab:compresion}
\centering
\small
\begin{tabular}{lrrrr}
\toprule
Dispositivo & $f$ (Hz) & $\Delta_C^{\text{mín}}$ (dB) & $\Delta_C^{\text{máx}}$ (dB) & $\delta$ (dB) \\
\midrule
Iron Deck  & 100  & $-3.25$ & $-0.02$ & $3.23$ \\
Iron Deck  & 1000 & $-3.26$ & $-0.00$ & $3.26$ \\
Iron Deck  & 5000 & $-3.25$ & $+0.01$ & $3.26$ \\
\midrule
Referencia & 100  & $-1.93$ & $-0.00$ & $1.93$ \\
Referencia & 1000 & $-1.44$ & $-0.00$ & $1.44$ \\
Referencia & 5000 & $-1.14$ & $+0.02$ & $1.16$ \\
\bottomrule
\end{tabular}
\end{table}

Los datos de la Tabla\,\ref{tab:compresion} revelan que Iron Deck exhibe una mayor excursión
de compresión ($\delta \approx 3.25\,\text{dB}$) uniforme en las tres bandas, consistente
con una curva de magnetización de tipo \emph{soft-knee} simétrica. En la referencia, la
no linealidad decrece con la frecuencia ($\delta = 1.93\,\text{dB}$ a $100\,\text{Hz}$,
$1.16\,\text{dB}$ a $5\,\text{kHz}$), documentando el efecto físico de pre-énfasis
de grabación a alta frecuencia que reduce la excursión de flujo y, por tanto, la saturación.

\subsection{THD, Anti-Aliasing y Distribución Espectral}
\label{sec:thd}

La Fig.\,\ref{fig:05_thd} grafica la THD (\%) a $1\,\text{kHz}$ en función del nivel de
entrada, dominada por armónicos impares ($H_3$, $H_5$) en Iron Deck, consistente con el
modelo de histéresis magnética de Preisach-Jiles-Atherton. La Fig.\,\ref{fig:06_aliasing}
demuestra la efectividad del filtro de sobremuestreo de Iron Deck: a $5\,\text{kHz}$ y
$0\,\text{dBFS}$, los picos armónicos se encuentran en los múltiplos exactos de la fundamental
sin energía espuria entre ellos, evidenciando la ausencia de \emph{foldover} hasta la
frecuencia de Nyquist.

\begin{figure}[htbp]
\centering
\includegraphics[width=0.48\textwidth]{resultados_analisis/05_THD_vs_Nivel.png}
\caption{Distorsión Armónica Total (THD\,\%) vs.\ nivel de entrada a $1\,\text{kHz}$,
$15\,\text{ips}$, NAB. Izquierda: Iron Deck (tres puntos de drive REC). Derecha:
dispositivo de referencia (tres niveles de Input).}
\label{fig:05_thd}
\end{figure}

\begin{figure}[htbp]
\centering
\includegraphics[width=0.48\textwidth]{resultados_analisis/06_Aliasing_5kHz_Saturacion_Extrema.png}
\caption{Prueba de anti-aliasing con tono de $5\,\text{kHz}$ a saturación máxima
($0\,\text{dBFS}$). Magnitud normalizada a la fundamental. Los marcadores verticales
señalan los múltiplos de $5\,\text{kHz}$. La ausencia de energía espuria entre armónicos
confirma el filtrado de sobremuestreo efectivo.}
\label{fig:06_aliasing}
\end{figure}

La Fig.\,\ref{fig:07_psd} presenta la distribución de energía por bandas de octava bajo
excitación de ruido rosa, con panel diferencial de escala propia para detectar diferencias
de décimas de dB que pasarían desapercibidas en curvas superpuestas.

\begin{figure}[htbp]
\centering
\includegraphics[width=0.48\textwidth]{resultados_analisis/07_Distribucion_Espectral_RuidoRosa.png}
\caption{Distribución de energía espectral por bandas de octava (ruido rosa, $-12\,\text{dBFS}$,
$15\,\text{ips}$, NAB). Panel izquierdo: Iron Deck vs.\ Referencia a ajuste nominal. Panel
derecho: Iron Deck a drive mínimo vs.\ máximo (migración espectral). El panel diferencial
inferior en cada caso aísla la variación banda a banda en dB.}
\label{fig:07_psd}
\end{figure}

\subsection{Estabilidad Digital de Frecuencia}
\label{sec:estabilidad}

La Tabla\,\ref{tab:estabilidad} resume la prueba de estabilidad digital mediante tonos patrón
de $3000\,\text{Hz}$ (NAB) y $3150\,\text{Hz}$ (DIN/IEC), procesados a todos los puntos de
drive disponibles. La frecuencia detectada se midió con Welch ($N_{\text{FFT}} = 4f_s$,
resolución $0.25\,\text{Hz}$).

\begin{table}[htbp]
\caption{Prueba de Estabilidad Digital de Frecuencia. La frecuencia detectada se midió
con resolución $0.25\,\text{Hz}$ (Welch, $N_\text{FFT}=4f_s$). Se reporta una fila
representativa por combinación de dispositivo y estándar; todos los puntos de drive
evaluados arrojaron $\text{Deriva}=+0.00\,\text{Hz}$.}
\label{tab:estabilidad}
\centering
\small
\begin{tabular}{llccc}
\toprule
Dispositivo & EQ & IPS & Patrón (Hz) & Deriva (Hz) \\
\midrule
Iron Deck & NAB & 7.5 & 3000 & $+0.00$ \\
Iron Deck & NAB & 15  & 3000 & $+0.00$ \\
Iron Deck & NAB & 30  & 3000 & $+0.00$ \\
Iron Deck & IEC & 30  & 3150 & $+0.00$ \\
\midrule
Referencia & NAB  & 7.5 & 3000 & $+0.00$ \\
Referencia & NAB  & 15  & 3000 & $+0.00$ \\
Referencia & CCIR & 15  & 3150 & $+0.00$ \\
\bottomrule
\end{tabular}
\end{table}

Ambos dispositivos demuestran deriva de frecuencia nula ($+0.00\,\text{Hz}$) bajo todos los
puntos de drive evaluados, confirmando que ni Iron Deck ni el dispositivo de referencia
introducen modulación de frecuencia ni artefactos de reloj digital.

% ===========================================================================
\section{Discusión}
% ===========================================================================

\begin{enumerate}
    \item \textbf{Respuesta espectral a $30\,\text{ips}$ (IEC/AES):} Iron Deck ofrece una
    extensión lineal hasta $22\,\text{kHz}$ que resulta ventajosa en flujos de trabajo a
    alta resolución, donde la atenuación de presencia en alta frecuencia introduciría
    coloración no deseada. La referencia, al limitarse a $15\,\text{ips}$ con CCIR, es
    más adecuada para emulación de estilo de producción europeo de los años 60--70.

    \item \textbf{Mantenimiento de calibración bajo saturación:} La invariancia $|\Delta H(f)|
    < 0.15\,\text{dB}$ de Iron Deck frente a incrementos de REC permite explotar la saturación
    armónica como herramienta de coloración sin comprometer el balance espectral de la mezcla.
    Este comportamiento es análogo al de una máquina física correctamente alineada con bias
    óptimo: la saturación no desplaza el punto de operación de la EQ de reproducción.

    \item \textbf{No linealidad dependiente de la frecuencia:} La excursión de compresión
    uniforme de Iron Deck ($\delta \approx 3.25\,\text{dB}$ en las tres bandas) documenta
    un modelo de histéresis simétrico. La reducción de $\delta$ con la frecuencia en la
    referencia ($3.0 \to 1.2\,\text{dB}$ de $100\,\text{Hz}$ a $5\,\text{kHz}$) es
    consistente con el efecto físico de pre-énfasis de grabación que reduce la excursión de
    flujo en alta frecuencia, disminuyendo la probabilidad de saturación.

    \item \textbf{Divergencia de estándares:} La elección entre IEC/AES y CCIR responde a
    filosofías de uso distintas. Ninguna norma es objetivamente superior; su pertinencia depende
    del repertorio, el periodo de grabación que se desea emular y el contexto de aplicación
    (mastering de alta resolución vs.\ coloración de estudio clásico europeo).

    \item \textbf{Integridad digital:} La deriva de frecuencia nula y la ausencia de aliasing
    espurio en ambos plugins confirman que las implementaciones digitales son coherentes con los
    requisitos de la norma AES17 \cite{aes17} para equipos de audio digital.
\end{enumerate}

% ===========================================================================
\section{Conclusión}
% ===========================================================================

La evaluación objetiva de emulaciones de cinta magnética mediante métricas DSP reproducibles
permite comparar dispositivos independientemente de su origen comercial. Los resultados
empíricos demuestran que Iron Deck (M~Media Audio) iguala la integridad digital del procesador
de referencia en las métricas críticas de estabilidad de frecuencia y anti-aliasing, y ofrece
ventajas técnicas en versatilidad de selección de velocidad ($7.5$\,/\,$15$\,/\,$30\,\text{ips}$)
y estándar de ecualización (NAB\,/\,IEC). El mantenimiento de calibración espectral bajo
saturación intensa ($|\Delta H(f)|<0.15\,\text{dB}$) es una ventaja operativa concreta para
el ingeniero de mezcla. Este trabajo establece un protocolo metodológico de \emph{benchmarking}
reproducible para la evaluación objetiva de emulaciones de cinta magnética en la comunidad de
ingeniería de audio.

\begin{thebibliography}{9}

\bibitem{aes17}
Audio Engineering Society, \emph{AES17-2020: AES standard method for digital audio engineering
--- Measurement of digital audio equipment}, New York, 2020.

\bibitem{nab}
National Association of Broadcasters, \emph{NAB Standard: Reproducing and Recording
Characteristics for Magnetic Tape Recording}, Washington, D.C., 1965.

\bibitem{iec94}
International Electrotechnical Commission, \emph{IEC 60094: Magnetic tape sound recording and
reproducing systems}, Geneva, Switzerland.

\bibitem{mrl}
Magnetic Reference Laboratory, \emph{MRL Calibration Tapes: Technical Notes on Fluxivity and
Equalization Standards}, San Jose, CA, 2019.

\bibitem{wallace}
R.\,L.\ Wallace, ``The Reproduction of Magnetically Recorded Signals,'' \emph{Bell System
Technical Journal}, vol.~30, no.~4, pp.~1145--1173, 1951.

\bibitem{zolzer}
U.\ Zölzer, \emph{Digital Audio Signal Processing}, 2nd ed., Chichester, UK: John Wiley \& Sons,
2008.

\end{thebibliography}

\end{document}
